\documentclass[12pt]{article}
\usepackage[spanish]{babel}
\usepackage[utf8]{inputenc}
\usepackage{amsmath, amssymb, graphicx}
\usepackage{geometry}
\geometry{margin=1in}

\title{CUANTIFICACIÓN DE LA ENERGÍA SOLAR PARA LA CIUDAD DE ORURO}
\author{M.Sc. Ing. Edgar S. Peñaranda Muñoz}
\date{}

\begin{document}
\maketitle

\section*{Resumen}
Luego de cinco años ininterrumpidos del servicio la estación meteorológica...
[RESUMEN COMPLETO TRANSCRITO]

\section*{Introducción}
Conformado el Instituto de Investigación...
[CONTENIDO COMPLETO]

\section*{Especificaciones geográficas y técnicas}
Latitud: $17^\circ 59' 36'' S$ \\
Longitud: $67^\circ 08' 09'' W$ \\
Altitud: 3725 m.s.n.m.

\section*{Parámetros solares}

\subsection*{Declinación solar}
\begin{equation}
\delta = 45.23 \sin\left(\frac{360}{365.25}(n-81)\right)
\end{equation}

\subsection*{Ángulo horario}
\begin{equation}
w = 15^\circ (H - 12)
\end{equation}

\subsection*{Altura solar}
\begin{equation}
\sin(ALT) = \cos(LAT)\cos(DEC)\cos(w) + \sin(LAT)\sin(DEC)
\end{equation}

\subsection*{Azimut solar}
\begin{equation}
\sin(AZ) = \frac{\sin(w)\cos(DEC)}{\cos(ALT)}
\end{equation}

\section*{Radiación solar diaria}

\begin{equation}
HHN = \sum IHN \cdot t
\end{equation}

\subsection*{Radiación extraterrestre}
\begin{equation}
H_0 = \frac{24}{\pi} I_0 \left[ \cos(LAT)\cos(DEC)\sin(WHE) + WHE\sin(LAT)\sin(DEC) \right]
\end{equation}

\subsection*{Índice de nubosidad}
\begin{equation}
K = \frac{HHN}{H_0}
\end{equation}

\subsection*{Fracción difusa (Liu \& Jordan)}
\begin{equation}
FD = 1.3903 - 4.0273K + 5.5315K^2 - 3.108K^3
\end{equation}

\subsection*{Componentes}
\begin{equation}
DFN = FD \cdot HHN
\end{equation}

\begin{equation}
DRN = HHN - DFN
\end{equation}

\section*{Radiación en superficies inclinadas}

\begin{equation}
RB = \frac{\int_{WE}^{WS} \cos(INC)\,dw}{\int_{WHS}^{WHE} \cos(INCH)\,dw}
\end{equation}

\begin{equation}
HDR = DRN \cdot RB
\end{equation}

\begin{equation}
HDF = DFN \cdot \frac{1 + \cos\gamma}{2}
\end{equation}

\begin{equation}
HRF = DFN \cdot RF \cdot \frac{1 - \cos\gamma}{2}
\end{equation}

\begin{equation}
HIT = HDR + HDF + HRF
\end{equation}

\section*{Radiación instantánea}

\begin{equation}
rd = \frac{IDFH}{DFN} = \frac{\pi}{24} \cdot \frac{\cos w - \cos WHE}{\sin WHE - WHE \cos WHE}
\end{equation}

\begin{equation}
IDRH = IHN - IDFH
\end{equation}

\section*{Superficies inclinadas (instantánea)}

\begin{equation}
rb = \frac{IDR}{IDRH} = \frac{\cos(INC)}{\cos(INCH)}
\end{equation}

\begin{equation}
IDF = IDFH \cdot \frac{1 + \cos\gamma}{2}
\end{equation}

\begin{equation}
IRF = IHN \cdot RF \cdot \frac{1 - \cos\gamma}{2}
\end{equation}

\begin{equation}
IT = IDR + IDF + IRF
\end{equation}

\section*{Resultados}

\subsection*{Tabla: Radiación diaria horizontal}
\begin{tabular}{lcccc}
\hline
 & Verano & Otoño & Invierno & Primavera \\
\hline
HHN & 6458 & 6326 & 5149 & 6785 \\
DRN & 4361 & 4564 & 4201 & 5345 \\
DFN & 2097 & 1761 & 948 & 1440 \\
\hline
\end{tabular}

\subsection*{Tabla: Superficie inclinada}
\begin{tabular}{lcccc}
\hline
 & Verano & Otoño & Invierno & Primavera \\
\hline
HIT & 5462 & 6377 & 6948 & 7163 \\
HDR & 3411 & 4648 & 6004 & 5737 \\
HDF & 1999 & 1679 & 904 & 1372 \\
HRF & 51 & 50 & 41 & 54 \\
\hline
\end{tabular}

\section*{Conclusiones}
La magnitud del histórico, su continuidad y calidad permiten cuantificar con fiabilidad la energía solar en Oruro...

\section*{Referencias}
Duffie, J.A., Beckman W.A., Solar Engineering, 1978 \\
Manrique, J.A., Energía Solar, 1984 \\
Cabrera, D., Sistemas Solares, 1996

\end{document}
